\section{Controlling the equality projections}
\label{sl:section}

We now bound the summed errors in Proposition~\ref{dis:interaction}.
On each pair fiber, equal-assignment indicators are quadratic functions
of the subset coordinates. Their span therefore uses only the
degree-one and degree-two parts of the Johnson slice decomposition;
see Filmus~\cite[Lemmas~18 and~24]{Filmus16}. We verify the required
identities directly. After the slice comparison, we return to the graph
space and sum the disjoint-pair inequalities.

\begin{lemma}[Equality-indicator projections on a slice]
\label{sl:comparison}
Let $N\ge0$ and $0\le q\le N$ be integers. Give
\[
 \mathcal X_{N,q}
 =\left\{x\in\{0,1\}^N:\ \sum_{i=1}^N x_i=q\right\}
\]
the uniform measure. Let $\tau_{ij}$ interchange coordinates $i,j$, and
put
\[
 G=\frac12\sum_{i<j}(I-\tau_{ij}).
\]
For each pair $i<j$, let $Q_{ij}$ be orthogonal projection onto the centered
function
\[
 z_{ij}=\mathbf 1_{\{x_i=x_j\}}
            -\mathbb P(x_i=x_j),
\]
with the convention $Q_{ij}=0$ when $z_{ij}=0$. Then
\begin{equation}
 S:=\sum_{i<j}Q_{ij}\preceq G.
 \label{sl:operator-comparison}
\end{equation}
\end{lemma}

\begin{proof}
Each $(I-\tau_{ij})/2$ is an orthogonal projection, so $G$ is positive
semidefinite. If all the $z_{ij}$ vanish, the assertion follows immediately.
This covers $N\le2$ and $q\in\{0,N\}$. In the remaining cases $N\ge3$ and
$1\le q\le N-1$. Writing $r=N-q$, we have
\begin{equation}
 \begin{aligned}
 p&:=\mathbb P(x_i=x_j)
   =\frac{q(q-1)+r(r-1)}{N(N-1)},\\
 1-p&=\frac{2qr}{N(N-1)},\qquad
 \|z_{ij}\|^2=p(1-p)>0.
 \end{aligned}
 \label{sl:p}
\end{equation}
All norms of functions in this proof use the uniform measure on the slice.

Let $M=\binom N2$ and give the pair coefficient space $\mathbb R^M$ its
usual Euclidean inner product. Define
\[
 Bc=\sum_{i<j}c_{ij}\frac{z_{ij}}{\sqrt{p(1-p)}}.
\]
Then $S=BB^*$. We analyze its Gram matrix $K=B^*B$ using the orthogonal
decomposition
\begin{align*}
 \mathcal C_0&=\{c:\ c_{ij}\text{ is constant}\},\\
 \mathcal C_1&=\left\{c:\ c_{ij}=u_i+u_j,\quad\sum_i u_i=0\right\},\\
 \mathcal C_2&=\left\{c:\ \sum_{j\ne i}c_{ij}=0\text{ for every }i\right\}.
\end{align*}
Here and below $c_{ji}=c_{ij}$. To verify the decomposition, the map
$u\mapsto(u_i+u_j)_{i<j}$ is injective for $N\ge3$: its kernel vanishes on
every triangle, and hence on every vertex. Its image is
$\mathcal C_0\oplus\mathcal C_1$, and its orthogonal complement is precisely
$\mathcal C_2$. Thus the dimensions are $1$, $N-1$, and $M-N$, respectively.
In particular $\mathcal C_2=\{0\}$ when $N=3$.

The number of equal coordinate pairs is fixed on the slice:
\[
 \sum_{i<j}\mathbf 1_{\{x_i=x_j\}}=\binom q2+\binom r2.
\]
It follows that $B$ kills $\mathcal C_0$. Permutation invariance shows that
the entries of $K$ depend only on the intersection size of their indexing
pairs. Moreover, for a fixed pair $ij$, the sums of coefficients over
\emph{other} pairs meeting $ij$, and over pairs disjoint from $ij$, are
\[
 \begin{array}{c|cc}
 &\text{intersection of size one}&\text{disjoint}\\ \hline
 c\in\mathcal C_1&(N-4)c_{ij}&-(N-3)c_{ij}\\
 c\in\mathcal C_2&-2c_{ij}&c_{ij}.
 \end{array}
\]
Indeed, the first column is the sum of the two row sums minus $2c_{ij}$;
the row sums are $(N-2)u_i$ on $\mathcal C_1$ and zero on $\mathcal C_2$.
The second column follows because both spaces have zero total coefficient
sum. For $N=3$ there are no disjoint pairs, and the displayed expression
on $\mathcal C_1$ is zero as required. Consequently $K$ acts by a scalar
$\mu_1$ on $\mathcal C_1$ and, when present, by a scalar $\mu_2$ on
$\mathcal C_2$.

We first compute $\mu_1$. The number of coordinates other than $i$ equal
to $x_i$ is
\[
 \sum_{j\ne i}\mathbf 1_{\{x_i=x_j\}}
   =(r-1)+(2q-N)x_i.
\]
For $\sum_i u_i=0$ this gives
\begin{equation}
 \sum_{i<j}(u_i+u_j)z_{ij}
       =(2q-N)\sum_i u_i x_i.
 \label{sl:linear-image}
\end{equation}
The one- and two-coordinate probabilities of a uniform $q$-subset imply
\[
 \operatorname{Var}\!\left(\sum_i u_i x_i\right)
    =\frac{qr}{N(N-1)}\sum_i u_i^2,
 \qquad
 \sum_{i<j}(u_i+u_j)^2=(N-2)\sum_i u_i^2.
\]
Using these identities and \eqref{sl:p} in
$\|Bc\|^2=\mu_1\|c\|^2$ yields
\begin{equation}
 \mu_1=\frac{(2q-N)^2}{2(N-2)p}\le\frac N2.
 \label{sl:mu-one}
\end{equation}
For the inequality, the difference of the relevant numerators is
\[
 N(N-2)p-(2q-N)^2
   =\frac{2N}{N-1}\bigl(qr-(N-1)\bigr)\ge0,
\]
since $q,r\ge1$ and $q+r=N$.

The diagonal entries of $K$ are all one, so $\operatorname{tr}K=M$.
Since $K$ is positive semidefinite, for $N\ge4$ we obtain
\begin{equation}
 (M-N)\mu_2\le M,
 \qquad
 0\le\mu_2\le\frac{N-1}{N-3}\le N-1.
 \label{sl:mu-two}
\end{equation}
This argument is used only when $\mathcal C_2$ is present.

To compare with $G$, set $\mathcal U_i=B\mathcal C_i$ for $i=1,2$.
These images are orthogonal, since
$\langle Bc,Bd\rangle=\langle c,Kd\rangle$. On $\mathcal U_i$ the operator
$S=BB^*$ acts by $\mu_i$; an image may of course be zero. Direct summation
of the transpositions gives
\begin{align}
 Gx_i&=\frac{Nx_i-q}{2},\label{sl:G-linear}\\
 G(x_ix_j)&=(N-1)x_ix_j-\frac{q-1}{2}(x_i+x_j).
 \label{sl:G-quadratic}
\end{align}
For the second identity, only swaps involving exactly one of $i,j$
contribute. Their sum is
\[
 (N-2)x_ix_j-\frac12(x_i+x_j)\sum_{k\notin\{i,j\}}x_k
 =(N-1)x_ix_j-\frac{q-1}{2}(x_i+x_j),
\]
using $x_i^2=x_i$ and $\sum_kx_k=q$.

By \eqref{sl:linear-image} and \eqref{sl:G-linear}, $G$ acts on
$\mathcal U_1$ by $N/2$. For $c\in\mathcal C_2$, the expansion
\[
 z_{ij}=1-p-x_i-x_j+2x_ix_j
\]
shows that $Bc$ is proportional to $\sum_{i<j}c_{ij}x_ix_j$: the constant
and linear terms vanish by the row-zero conditions. Those same conditions
cancel the linear terms in \eqref{sl:G-quadratic}, so $G$ acts on
$\mathcal U_2$ by $N-1$. Equations \eqref{sl:mu-one} and
\eqref{sl:mu-two} therefore give $S\preceq G$ on
$\mathcal U_1\oplus\mathcal U_2$. This space is invariant under the
self-adjoint $G$, hence so is its orthogonal complement. On that complement
$S=0$, because $B\mathcal C_0=0$ and the three coefficient spaces exhaust
$\mathbb R^M$; meanwhile $G\succeq0$. This proves
\eqref{sl:operator-comparison} on the whole slice.

For clarity, the possible zero images introduce no additional cases:
when $q=1$ or $q=N-1$, \eqref{sl:mu-one} gives $\mu_1=N/2$, and the trace
identity gives $\mu_2=0$ if $N\ge4$; when $q=N/2$, $\mu_1=0$.
For $N=3$ the second image is absent. At $N=4,q=2$, the trace identity gives
$\mu_2=3=N-1$, so the estimate includes this boundary case with equality.
\end{proof}

\begin{corollary}[Summed equality-indicator energy]
\label{sl:global}
For every vertex pair $a$ and every function $f$ on the graph space,
\begin{equation}
 \sum_{b:\,b\cap a=\varnothing}\|R_{a,b}f\|^2
 \le
 \sum_{b:\,b\cap a=\varnothing}\|g_{a,b}f\|^2.
 \label{sl:graph-comparison}
\end{equation}
\end{corollary}

\begin{proof}
Work first in one $a$-fiber, whose assignments form a uniform $q$-subset
slice of its $N$ singleton vertices. If both vertices of $b$ are singleton
labels $k,l$, the involution $T_{a,b}$ acts on this slice exactly as
$\tau_{kl}$: it interchanges their assignments when they differ and acts
trivially when they agree. Thus $g_{a,b}=(I-\tau_{kl})/2$, while
$R_{a,b}=Q_{kl}$ in the notation of Lemma~\ref{sl:comparison}.
If either vertex of $b$ is not a singleton, the cross matrix cannot have
both column sums equal to one, and both $g_{a,b}$ and $R_{a,b}$ vanish
on the fiber. Since an orthogonal projection $Q$ satisfies
$\|Qf\|^2=\langle f,Qf\rangle$, the two conditional sums in
\eqref{sl:graph-comparison} are $\langle f,Sf\rangle$ and
$\langle f,Gf\rangle$. Lemma~\ref{sl:comparison} compares them.
Both operators preserve the complete $a$-fibers, so averaging these
conditional inequalities proves the corollary.
\end{proof}

\begin{corollary}\label{dis:nonnegative}
The ordered sum
\begin{equation}\label{dis:D}
 D=\sum_{\substack{a,b\in\binom V2\\a\cap b=\varnothing}}h_a h_b
\end{equation}
is self-adjoint and positive semidefinite.
\end{corollary}

\begin{proof}
The two orders of every disjoint pair occur in the sum, proving
self-adjointness. Grouping them and applying
Proposition~\ref{dis:interaction} gives
\begin{align*}
 \langle f,Df\rangle
 &\ge\sum_{\{a,b\}:\,a\cap b=\varnothing}
       \bigl(2\|g_{a,b}f\|^2
             -\|R_{a,b}f\|^2-\|R_{b,a}f\|^2\bigr)\\
 &=\sum_{a\in\binom V2}
       \left(\sum_{b:\,a\cap b=\varnothing}\|g_{a,b}f\|^2
            -\sum_{b:\,a\cap b=\varnothing}\|R_{a,b}f\|^2\right)
 \ge0.
\end{align*}
The first sum is over unordered pairs of disjoint vertex pairs;
the equality uses $g_{a,b}=g_{b,a}$. The last inequality is
Corollary~\ref{sl:global}, applied separately for each $a$.
\end{proof}
